\addcontentsline{toc}{section}{СПИСОК ИСПОЛЬЗОВАННЫХ ИСТОЧНИКОВ}

\begin{thebibliography}{9}
    
    \bibitem{1} Бретт, М. PHP и MySQL. Исчерпывающее руководство / М. Бретт. – Санкт-Петербург : Питер, 2016. – 544 с. – ISBN 978-5-496-01049-8. – Текст~: непосредственный.
    \bibitem{2} Веру, Л. Секреты CSS. Идеальные решения ежедневных задач / Л. Веру. – Санкт-Петербург : Питер, 2016. – 336 с. – ISBN 978-5-496-02082-4. – Текст~: непосредственный.
	\bibitem{3}	Титтел, Э. HTML5 и CSS3 для чайников / Э. Титтел, К. Минник. – Москва~: Вильямс, 2016 – 400 с. – ISBN 978-1-118-65720-1. – Текст~: непосредственный.    
	\bibitem{4}	Макнейл П. Веб-дизайн. Книга идей веб-разработчика / П. Макнейл. — СПб.: Питер, 2017. — 480 c. – ISBN 978-5-496-00705-4. – Текст~: непосредственный.
    \bibitem{5} Леонова В. Н. Основы продвижения сайтов : методические указания / В. Н. Леонова. — Гомель : Гомельский государственный технический университет имени П. О. Сухого, 2026. — Текст~: непосредственный.
    \bibitem{6} Макфедрис П. Веб-дизайн с нуля: HTML + CSS на практике / П. Макфедрис ; пер. с англ. — 2-е изд. — Санкт-Петербург : Питер, 2025. — 416 c. — ISBN 978-5-4461-4128-9. — Текст~: непосредственный.
    \bibitem{7} Бачурина Т. В. Веб-дизайн и реклама / Т. В. Бачурина, О. М. Кунцевич. — Минск : Белорусский государственный университет культуры и искусств, 2025. — Текст~: непосредственный.
    \bibitem{8}	 Нильсен Я. Веб-дизайн: книга Якоба Нильсена / Я. Нильсен. — М.: Символ, 2015. — 512 c. – ISBN 978-5-93286-004-5. – Текст~: непосредственный.
    \bibitem{9} Петин В. А. Веб-дизайн и юзабилити. Практическое руководство / В. А. Петин. — 2-е изд., перераб. и доп. — Москва : ДМК Пресс, 2025. — 352 c. — ISBN 978-5-93700-412-1. — Текст~: непосредственный.
    \bibitem{10} Морозов Д. А. Frontend-разработка. Реактивный подход и современные фреймворки / Д. А. Морозов. — Москва : Солон-Пресс, 2025. — 304 c. — ISBN 978-5-91359-563-8. — Текст~: непосредственный.
    \bibitem{11} Гурли,Д.HTTP.Подробноеруководство/Д.Гурли,Б.Тотти.–Санкт-Петербург : Символ-Плюс, 2013.– 656 с.– ISBN 978-5-93286-209-7.– Текст : непосредственный.
    \bibitem{12} Макфарланд, Д. Большая книга CSS / Д. Макфарланд.– Санкт-Петербург : Питер, 2012.– 560 с.– ISBN 978-5-496-02080-0.– Текст : непосредственный.
    \bibitem{13} Немцова Т. И., Казанкова Т. В., Шнякин А. В. Компьютерная графика и web-дизайн. Учебное пособие. — М.: Форум, 2023. — 400 c.– ISBN 978-5-16-021098-8. Текст : непосредственный.
    \bibitem{14} Голдстайн, А. HTML5иCSS3длявсех/А.Голдстайн, Л. Лазарис, Э.Уэйл.– Москва : Вильямс, 2012.– 368 с.– ISBN 978-5-699-57580-0.– Текст :непосредственный.
    \bibitem{15} Дэкетт, Д. HTML и CSS. Разработка и создание веб-сайтов / Д. Дэкетт.– Москва : Эксмо, 2014.– 480 с.– ISBN 978-5-699-64193-2.– Текст :непосредственный.
    \bibitem{16} Хавербеке, М. Выразительный JavaScript / М. Хавербеке.– Санкт-Петербург : Питер, 2019.– 472 с.– ISBN 978-5-4461-1211-1.– Текст : непосредственный.
    \bibitem{17} Крокфорд, Д. JavaScript: сильные стороны / Д. Крокфорд.– Санкт-Петербург : Питер, 2012.– 176 с.– ISBN 978-5-459-00788-6.– Текст : непосредственный.
    \bibitem{18} Закас Н.С. Основы JavaScript для профессиональных веб-разработчиков / Н. Закас; пер. с англ. — 4-е изд. — Санкт-Петербург: Вильямс, 2020. — 1216 с. — ISBN 978-5-8459-2394-9.
    \bibitem{19} Флэнаган Д. JavaScript. Подробное руководство / Д. Флэнаган; пер. с англ. — 7-е изд. — Санкт-Петербург: БХВ-Петербург, 2021. — 1168 с. — ISBN 978-5-9775-4825-9.
    \bibitem{20} Прохоренок Н.А., Дронов В.А. JavaScript и Node.js для веб-разработчиков / Н.А. Прохоренок, В.А. Дронов. — Санкт-Петербург: БХВ Петербург, 2022. — 592 с. — ISBN 978-5-9775-6591-9.
    \bibitem{21} Макфарланд Д. Изучаем HTML, XHTML и CSS / Д. Макфарланд; пер. с англ. — 2-е изд. — Санкт-Петербург: Питер, 2020. — 656 с. — ISBN 978-5-4461-1244-8.
    \bibitem{22} Выдрин Д.К., Грибов А.Ю., Анохин А.Ю. HTML5 и CSS3. Веб-разработка по стандартам нового поколения / Д.К. Выдрин, А.Ю. Грибов, А.Ю. Анохин. — 2-е изд., перераб. и доп. — Москва: ДМК Пресс, 2021. — 416 с. — ISBN 978-5-97060-957-9.
    \bibitem{zakon1} Российская Федерация. Законы. О защите прав потребителей: Закон РФ от 7 фервраля 1992г N 2300-1 – Текст : электронный // КонсультантПлюс: справочно-правовая система. URL:https://www.consultant.ru/document/cons\_doc\_LAW\_305/ (дата обращения: 04.05.2026).
    \bibitem{zakon2} Российская Федерация. Законы. Глава 34 «Аренда». Гражданский кодекс РФ – Текст : электронный // КонсультантПлюс: справочно-правовая система. — URL: https://www.consultant.ru/document/cons\_doc\_LAW\_9027
    /cac65a05697011caaa0eff7d476d867137f1f92a/ (дата обращения: 04.05.2026).
    \bibitem{zakon3} Российская Федерация. Законы. О проведении эксперимента по установлению специального налогового режима "Налог на профессиональный доход": Федеральный закон от 27.11.2018 № 422-ФЗ – Текст : электронный // КонсультантПлюс: справочно-правовая система. URL:https://www.consultant.ru/document/cons\_doc\_LAW\_311977/ (дата обращения: 04.05.2026).
    \bibitem{zakon4} Российская Федерация. Законы. О персональных данных: Федеральный закон от 27.07.2006 № 152-ФЗ – Текст : электронный // КонсультантПлюс: справочно-правовая система. URL:https://www.consultant.ru/document/cons\_doc\_LAW\_61801/ (дата обращения: 05.05.2026).
    \bibitem{gost701} ГОСТ 19.701-90. Единая система программной документации. Схемы алгоритмов, программ, данных и систем.– Москва : Стандартинформ, 2010. – Текст~: непосредственный.
    \bibitem{gost602} ГОСТ 34.602-2020. Информационные технологии. Комплекс стандартов на автоматизированные системы. Техническое задание на создание автоматизированной системы.– Москва : Стандартинформ, 2020. – Текст~: непосредственный.
    \bibitem{gost201} ГОСТ 34.201-2020. Информационные технологии. Комплекс стандартов на автоматизированные системы. Виды,комплектность и обозначение документов при создании автоматизированных систем.– Москва : Стандарт информ, 2020. – Текст~: непосредственный.
    \bibitem{23} Алексеев А. П. Введение в Web-дизайн. Учебное пособие. — М.: Солон-Пресс, 2021. — 184 c. – ISBN 978-5-91359-355-9. Текст~: непосредственный.
    \bibitem{24} Полуэктова Н. Р. Разработка веб-приложений. — М.: Юрайт, 2024. — 205 c. – ISBN 978-5-534-18644-4. Текст~: непосредственный.
    \bibitem{25} Петроченков А., Новиков Е. Идеальный Landing Page. Создаем продающие веб-страницы. — СПб.: Питер, 2017. — 320 c. – ISBN 978-5-4461-0292-1. Текст~: непосредственный.
    \bibitem{26} Кириченко А. В. Справочник HTML. Кратко, быстро, под рукой. — М.: Наука и техника, 2023. — 288 c. – ISBN 978-5-94387-275-4. Текст~: непосредственный.
    \bibitem{27} Хрусталев А. А. Дубовик Е. В. Справочник CSS3. Кратко, быстро, под рукой. — М.: Наука и техника, 2021. — 304 c. – ISBN 978-5-94387-279-2. Текст~: непосредственный.
    \bibitem{28} Кангин В. В. Интернет. Языки HTML и JavaScript. — М.: ТНТ, 2021. — 488 c. – ISBN 978-5-94178-140-9. Текст~: непосредственный.
    \bibitem{29} Стефанов С. React. Быстрый старт, 2-е изд. / С. Стефанов. — Санкт-Петербург : Питер, 2023. — 304 c. – ISBN 978-5-4461-2115-1. – Текст~: непосредственный.
    \bibitem{30} Заяц А. М. Проектирование и разработка WEB-приложений. Введение в frontend и backend разработку на JavaScript и node.js : учебное пособие для СПО / А. М. Заяц, Н. П. Васильев. — 3-е изд., стер. — Санкт-Петербург : Лань, 2023. — 120 c. – ISBN 978-5-507-46023-4. – Текст~: непосредственный.
    \bibitem{31} Сысолетин Е. Г. Разработка интернет-приложений : учебное пособие для СПО / Е. Г. Сысолетин, С. Д. Ростунцев, Л. Г. Доросинский. — Москва : Юрайт, 2024. — 80 c. – ISBN 978-5-534-19603-0. – Текст~: непосредственный.
    \bibitem{32} Коллер Ф. М. Полный стек. Веб-разработка / Ф. М. Коллер ; пер. с англ. — Москва : Эксмо, 2024. — 432 c. — ISBN 978-5-04-192845-6. — Текст~: непосредственный.
    \bibitem{33} Маллах Э. Введение в веб-разработку с HTML, CSS и JavaScript / Э. Маллах ; пер. с англ. — Москва : ДМК Пресс, 2025. — 414 c. — ISBN 978-5-93700-312-4. — Текст~: непосредственный.
    \bibitem{34} Никсон Р. HTML5 и CSS3. Мастер-класс / Р. Никсон ; пер. с англ. — Санкт-Петербург : БХВ-Петербург, 2024. — 464 c. — ISBN 978-5-9775-1929-8. — Текст~: непосредственный.
    \bibitem{35} Янцев В. В. Разработка web-страниц на HTML, CSS и JavaScript : учебное пособие для вузов / В. В. Янцев. — Санкт-Петербург : Лань, 2024. — 148 c. — ISBN 978-5-507-49407-1. — Текст~: непосредственный.
    \bibitem{36} Тузовский А. Ф. Проектирование и разработка web-приложений : учебник для вузов / А. Ф. Тузовский. — Москва : Юрайт, 2025. — 219 c. — ISBN 978-5-534-16300-1. — Текст~: непосредственный.
    \bibitem{37} Баркович А. А. Веб-проектирование : учебное пособие / А. А. Баркович, Т. А. Филимонова. — Москва : ИНФРА-М, 2025. — 231 c. — ISBN 978-5-16-019399-1. — Текст~: непосредственный.
    \bibitem{38} Фролов А. Б. Основы web-дизайна. Разработка, создание и сопровождение web-сайтов : учебное пособие для СПО / А. Б. Фролов, И. А. Нагаева, И. А. Кузнецов. — Саратов : Профобразование, 2024. — 144 c. — ISBN 978-5-4488-2231-5. — Текст~: непосредственный.
    \bibitem{39} Жемчужников Д. Г. Веб-дизайн. Уровень 2 : учебное пособие / Д. Г. Жемчужников. — Москва : Просвещение, 2025. — 112 c. — ISBN 978-5-09-087118-1. — Текст~: непосредственный.
    \bibitem{40} Кузнецова Л. В. Современные веб-технологии : учебное пособие / Л. В. Кузнецова. — Москва : Интернет-Университет Информационных Технологий, 2024. — 187 c. — ISBN 978-5-4497-2457-1. — Текст~: непосредственный.
    \bibitem{41} Булгакова И. А. Разработка и прототипирование веб-сайтов и интерфейсов онлайн : учебное пособие для вузов / И. А. Булгакова. — Москва : Дашков и К°, 2024. — 215 c. — ISBN 978-5-394-06214-8. — Текст~: непосредственный.
    \bibitem{42} Илларионова М. В. WEB-дизайн : учебное пособие / М. В. Илларионова. — Санкт-Петербург : Санкт-Петербургский политехнический университет Петра Великого, 2024. — 140 c. — DOI 10.18720/SPBPU/5/tr24-98. — Текст~: непосредственный.
    \bibitem{43} Медведев М. А. Web технологии в бизнесе : учебное пособие / М. А. Медведев, М. А. Медведева ; под ред. D. B. Berg. — Екатеринбург : Издательство Уральского университета, 2024. — 160 c. — ISBN 978-5-7996-3906-8. — Текст~: непосредственный. 
\end{thebibliography}
